\section{从精确插值到整体逼近}
\label{sec:08-Numerical-Analysis/03-Interpolation-and-Approximation/03}

前两节的所有讨论都默认了一件事：数据是对的。误差公式、Chebyshev 节点、Lebesgue 常数，全都建立在``\(y_i\) 就是 \(f(x_i)\)''这个前提上。现在把前提抽掉。

假设你手上是一批传感器读数，每个数带着几个百分点的随机偏移。你仍然可以做插值，多项式仍然会精确穿过每一个点——包括穿过每一次随机抖动。曲线弯下去不是因为函数在那里弯，是因为那次采样恰好偏低。

\begin{center}
\begin{tikzpicture}[x=1.32cm,y=1.02cm,font=\small,>=stealth]
  \draw[->,black!55] (-0.15,0) -- (6.5,0) node[right,font=\footnotesize] {\(x\)};
  \draw[->,black!55] (0,-0.15) -- (0,4.5) node[above,font=\footnotesize] {\(y\)};
  \foreach \t in {0,2,4,6} \node[below,font=\footnotesize] at (\t,-0.12) {\t};
  \draw[red!70,thick] plot[smooth] coordinates {
    (0.00,0.600) (0.20,1.850) (0.40,2.266) (0.60,2.197) (0.80,1.900) (1.00,1.550)
    (1.20,1.257) (1.40,1.080) (1.60,1.037) (1.80,1.120) (2.00,1.300) (2.20,1.540)
    (2.40,1.799) (2.60,2.039) (2.80,2.229) (3.00,2.350) (3.20,2.394) (3.40,2.367)
    (3.60,2.288) (3.80,2.187) (4.00,2.100) (4.20,2.069) (4.40,2.130) (4.60,2.312)
    (4.80,2.625) (5.00,3.050) (5.20,3.531) (5.40,3.960) (5.60,4.165) (5.80,3.894)
    (6.00,2.800)};
  \draw[blue!75,very thick] (0,0.850) -- (6,3.079);
  \foreach \p in {(0,0.60),(1,1.55),(2,1.30),(3,2.35),(4,2.10),(5,3.05),(6,2.80)}
    \fill[black!85] \p circle (1.7pt);
  \draw[black!50,dashed] (1,1.221) -- (1,1.55);
  \draw[black!50,dashed] (2,1.593) -- (2,1.30);
  \draw[black!50,dashed] (3,1.964) -- (3,2.35);
  \draw[black!50,dashed] (5,2.707) -- (5,3.05);
  \node[right,font=\footnotesize,red!75] at (4.4,4.35) {六次插值};
  \node[right,font=\footnotesize,blue!80] at (2.25,0.72) {最小二乘直线};
\end{tikzpicture}

\emph{七个带噪读数。红线穿过每一个点，于是在两点之间上下翻腾，右端还甩出一个高峰；蓝线一个点都不穿过（虚线段是残差），却更像那条产生数据的规律。数据可信时红线是答案，数据带噪时它是灾难。}
\end{center}

有人会说：那就少取几个点，用低次多项式插值。这确实能压住振荡，但被丢掉的点里也有信号，你是在用扔信息换稳定。真正该动的不是节点数，而是那条一直没被质疑的约束——\textbf{凭什么一定要穿过每个点}。松开它，问题就从插值变成逼近：不再要求逐点命中，只要求整体上离得近。

\subsection{不先说清用什么尺子，``最佳''就没有定义}

逼近问题的形式很短：在 \(n\) 次多项式全体 \(P_n\) 里找一个 \(p_n\)，使 \(\norm{f-p_n}\) 尽量小。麻烦全在那个范数符号上。

用 \(L^2\) 范数 \(\bigl(\int_a^b\abs{f-p_n}^2w\,dx\bigr)^{1/2}\)，你度量的是误差能量在整个区间上的加权平均。平方会放大大误差，但单个点上的一次剧烈偏离仍可能被积分摊薄。数据是离散的时候它退化成熟悉的残差平方和 \(\sum_i(y_i-p_n(x_i))^2\)。换成 \(L^\infty\) 范数 \(\max_{x}\abs{f(x)-p_n(x)}\)，你只盯着最差的那一个点，别处再好也不加分。再往里加权重函数，就是在声明区间的哪些部分你更在乎。

同一个 \(f\)、同一个 \(P_n\)，尺子一换，最优解就换。上图那七个点若用直线拟合，最小二乘给出 \(y=0.850+0.371x\)，残差平方和 \(0.656\)、最大偏差 \(0.386\)；一致逼近给出 \(y=0.897+0.371x\)，最大偏差降到 \(0.339\)，残差平方和却涨到 \(0.671\)。两条线各在自己的赛道上夺冠，彼此都不是对方眼里的最优。

这件事在训练模型时天天发生，只是换了个说法：用平方损失还是用最大误差，取决于你怕哪种错误。推荐系统里少数用户体验极差可以被平均掉，那就用平方损失；控制器的位置误差一旦越限就是事故，那就得盯最坏点。选损失函数不是调参，是在声明代价结构。

\subsection{\texorpdfstring{\(L^2\) 逼近就是一次正交投影}{L2 逼近就是一次正交投影}}

在 \(L^2\) 这把尺子下，逼近问题有个干净的几何解释。带上内积 \(\ip{f}{g}=\int_a^b f g\,w\,dx\)，``离 \(f\) 最近的 \(p_n\)''就是 \(f\) 在子空间 \(\operatorname{span}\set{\phi_0,\dots,\phi_n}\) 上的正交投影，判据是残差与子空间正交：\(\ip{f-p_n}{\phi_j}=0\) 对每个 \(j\) 成立。

这套说法你在\hyperref[sec:03-Linear-Algebra/03-Orthogonality/02]{``投影与最近点：从一条直线走向一般子空间''}里已经见过，那里的向量在 \(\R^n\) 里，这里的向量是函数。骨架完全一样：把 \(p_n=\sum_k c_k\phi_k\) 代进正交条件，若基两两正交，交叉项全灭，只剩

\[
c_j=\frac{\ip{f}{\phi_j}}{\ip{\phi_j}{\phi_j}}.
\]

联立求解 \(n+1\) 个系数的方程组，被拆成了 \(n+1\) 次独立的标量除法。这是正交基唯一的、也是全部的好处。

顺带看一个小例子，它能治好``逼近就是要贴着函数走''的错觉。在 \([-1,1]\) 上把 \(f(x)=x^2\) 投影到 \(\operatorname{span}\set{1,x}\)（取 \(w\equiv1\)）：一次项系数为零（偶函数与奇函数正交），常数项 \(c_0=\int_{-1}^1x^2dx\big/\int_{-1}^11\,dx=1/3\)。最佳 \(L^2\) 逼近是常函数 \(1/3\)，它与 \(x^2\) 在任何一个点上都不相等，\(x=0\) 处差 \(-1/3\)，\(x=1\) 处差 \(+2/3\)。它不讨好任何单点，只让平方误差的总账最小。

正交基的必要性还有另一面。直接用单项式 \(1,x,x^2,\dots\)，在有限区间上它们彼此太像了——\(x^{10}\) 与 \(x^{11}\) 在 \([0,1]\) 上几乎共线——正规方程的系数矩阵变成 Hilbert 型矩阵，条件数随次数指数上升，高次系数完全不可信。Legendre、Chebyshev 这些正交族就是为绕开这件事而生的。

离散最小二乘那边有个配套的实现纪律：给定 \(m\) 个点用 \(n\) 次多项式拟合（\(m>n\)），应对设计矩阵直接做 QR 或 SVD，而不是先拼出 \(A^{\trans}A\) 再解。正规方程会把条件数近似平方，本来只是勉强的问题会被推进灾难区。代码里冒出负方差或系数爆炸时，第一个该查的就是这里，详见\hyperref[sec:08-Numerical-Analysis/05-Numerical-Linear-Algebra/02]{``QR 与最小二乘的稳定解法''}。

\subsection{一致逼近：与最坏的那个点较劲}

\(L^\infty\) 逼近要解的是

\[
\min_{p\in P_n}\ \max_{x\in[a,b]}\abs{f(x)-p(x)}.
\]

它的最优解有一个可以直接用来验货的特征：误差函数 \(f-p^*\) 会在至少 \(n+2\) 个点上交替取到等幅的正负极值。

\begin{center}
\begin{tikzpicture}[x=1.15cm,y=1.05cm,font=\small,>=stealth]
  \draw[->,black!55] (-0.2,0) -- (7.1,0) node[right,font=\footnotesize] {\(x\)};
  \draw[black!30,dashed] (-0.1,1) -- (6.9,1);
  \draw[black!30,dashed] (-0.1,-1) -- (6.9,-1);
  \draw[blue!75,thick] plot[smooth] coordinates {
    (0.00,-1.00) (0.30,-0.86) (0.60,-0.50) (0.90,0.00) (1.20,0.52) (1.50,0.88)
    (1.70,1.00) (1.95,0.90) (2.25,0.55) (2.55,0.05) (2.85,-0.48) (3.15,-0.85)
    (3.40,-1.00) (3.70,-0.88) (4.00,-0.52) (4.30,0.00) (4.60,0.52) (4.90,0.88)
    (5.15,1.00) (5.45,0.88) (5.75,0.50) (6.05,-0.03) (6.35,-0.58) (6.65,-0.90)
    (6.90,-1.00)};
  \foreach \p in {(0,-1),(1.7,1),(3.4,-1),(5.15,1),(6.9,-1)}
    \fill[red!80] \p circle (1.7pt);
  \node[left,font=\footnotesize] at (-0.15,1) {\(+E\)};
  \node[left,font=\footnotesize] at (-0.15,-1) {\(-E\)};
  \node[right,font=\footnotesize,blue!75] at (2.4,1.35) {误差函数 \(f-p^{*}\)};
\end{tikzpicture}

\emph{最佳一致逼近的误差交替触碰 \(\pm E\)（红点），一次都不越界。若某个正峰明显低于别的峰，说明还有余量：微调系数把最高的几个峰压下去、把矮的抬上来，最大值还能再降。峰全部等高时，任何调整都会在别处顶出新的最高点，这就是最优的标志。}
\end{center}

Chebyshev 多项式和这个特征是同一件事的两面：\(T_n(x)=\cos(n\arccos x)\) 的极值恰好在 \(\pm1\) 之间交替，它正是``给定最高次系数，让最大绝对值最小''这一 minimax 问题的解。上一节把节点推向两端的做法，源头就在这里——Chebyshev 节点做的事，是让误差公式里那个节点乘积在一致范数下取到最优。

一致逼近的典型用户是数学函数库。实现 \(\exp\) 或 \(\log\) 时，厂商承诺的是全定义域内每一位有效数字都可信，不是``平均而言很准''。这种承诺只能由 \(L^\infty\) 兑现。

\subsection{Weierstrass 说存在，但只说到这里}

\begin{theorem}
设 \(f\) 在闭区间 \([a,b]\) 上连续。对任意 \(\varepsilon>0\)，存在多项式 \(p\) 使 \(\norm{f-p}_\infty<\varepsilon\)。
\end{theorem}

多项式在连续函数空间里稠密，这是逼近论的地基。证明有好几条路（Bernstein 多项式的概率证法最短），这里不展开，因为对使用者而言更重要的是它\emph{没有}承诺的东西。

它不说需要多高的次数。存在某个 \(n\)，但那个 \(n\) 可能大得没有工程意义，而且强烈依赖 \(f\) 的光滑程度：解析函数常有谱收敛，Lipschitz 连续但不可导的函数最多只能指望 \(O(1/n)\) 的阶。它也不说怎么把那个多项式找出来——上一节已经演示过，用等距节点插值这条最自然的路会通向 Runge 现象，存在性与可构造性之间隔着一整个方法论。它更不管浮点：定理站在实数算术上，而高次多项式的求值本身就是病态运算。

边界还有一条硬的：\(f\) 在区间内有跳跃时，一致逼近到任意小是不可能的。多项式连续，一致极限保持连续，跳点两侧的值没法被同一个连续函数同时追上。这时要么把跳点的小邻域挖掉，要么换范数——\(L^2\) 下单点的误差在积分里权重为零，结论就完全不同。范数选择再一次决定了什么叫做得到。

\subsection{插值还是逼近，取决于数据不是取决于品味}

两者不是精粗之分。插值把``穿过数据''当成硬约束，逼近把它换成软代价，这是不同的问题设定，不是同一问题的两种解法。

决定用哪个，问几件事就够了：数据是精确求值还是含噪测量；节点处残差为零是物理要求（比如边界条件必须严格满足）还是可以放松；你更在乎平均行为还是最坏情况；节点位置能不能自己挑；模型是用在数据范围之内还是要往外推。函数已知、无噪、光滑时，Chebyshev 节点上的高次插值往往就是最优解；测量含噪、样本量远大于自由度时，低次最小二乘更安全。选错的通常不是方法，是没先问清自己要最小化什么。

不过逼近也没有解决全部问题。最小二乘直线在整个区间上只有两个自由度，遇到形状复杂的曲线仍然力不从心；而提高多项式次数，Runge 那一套麻烦又会找回来。下一节走第三条路：既不用一条高次曲线统管全区间，也不放弃穿过数据——把区间切开，每段只用三次。
